\documentclass[a4paper,12pt]{article}

% --- Пакеты (XeLaTeX + кириллица) ---
\usepackage{fontspec}
\setmainfont{Times New Roman}
\usepackage[russian]{babel}
\usepackage{geometry} % Для настройки полей
\usepackage{titlesec} % Для форматирования заголовков
\usepackage{graphicx} % Для вставки изображений
\usepackage{booktabs} % Для красивых таблиц
\usepackage{hyperref} % Для ссылок
\usepackage{enumitem} % Для настройки списков

% --- Настройка страницы ---
\geometry{left=2.5cm}
\geometry{right=2.5cm}
\geometry{top=2.5cm}
\geometry{bottom=2.5cm}

% --- Настройка заголовков ---
\titleformat{\section}{\Large\bfseries}{\thesection}{1em}{}
\titleformat{\subsection}{\large\bfseries}{\thesubsection}{1em}{}

% --- Данные документа ---
\title{\textbf{БИЗНЕС-ПЛАН}\\[0.5em] \large Название вашего проекта}
\author{Автор: Имя Фамилия \\ Компания: Название компании}
\date{Дата: \today}

\begin{document}

% --- Титульная страница ---
\maketitle
\thispagestyle{empty}
\newpage

% --- Содержание ---
\tableofcontents
\newpage

% ==================== РАЗДЕЛ 1: РЕЗЮМЕ ====================
\section{Резюме}

Резюме — это краткое изложение всего бизнес-плана. Оно должно быть написано в последнюю очередь, но располагается в начале документа. Здесь нужно описать суть проекта, его цели и основные преимущества.

\subsection{Цель проекта}
Кратко опишите, чего вы хотите достичь (например: захват 10\% рынка за 2 года, выход на окупаемость через 18 месяцев).

\subsection{Основные показатели}
\begin{itemize}
    \item Объем необходимых инвестиций: \textbf{X рублей}.
    \item Срок окупаемости (PP): \textbf{Y месяцев}.
    \item Чистая приведенная стоимость (NPV): \textbf{Z рублей}.
\end{itemize}

% ==================== РАЗДЕЛ 2: ОПИСАНИЕ КОМПАНИИ ====================
\newpage
\section{Описание компании}

\subsection{О компании}
Опишите историю создания, юридический статус, местоположение и миссию вашей компании.

\begin{quote}
    \textit{«Наша миссия — сделать жизнь наших клиентов лучше благодаря...»}
\end{quote}

\subsection{Команда}
Представьте ключевых участников команды. Инвесторы инвестируют в людей, а не только в идеи.

\begin{itemize}
    \item \textbf{Генеральный директор:} Опыт работы, ключевые навыки.
    \item \textbf{Технический директор:} Опыт работы, ключевые навыки.
\end{itemize}

% ==================== РАЗДЕЛ 3: АНАЛИЗ РЫНКА ====================
\newpage
\section{Анализ рынка}

\subsection{Обзор отрасли}
Опишите текущее состояние отрасли, тенденции роста и прогнозы развития. Используйте статистику и аналитические данные.

\subsection{Целевая аудитория}
Кто ваш клиент? Опишите портрет вашего идеального покупателя.
\begin{itemize}
    \item Возраст, пол, доход.
    \item География.
    \item Потребности и боли клиента.
\end{itemize}

\subsection{Конкурентный анализ}
Сравните ваш проект с основными конкурентами.

\begin{table}[h!]
\centering
\renewcommand{\arraystretch}{1.5}
\begin{tabular}{|p{3cm}|p{3cm}|p{3cm}|p{3cm}|}
\hline
\textbf{Критерий} & \textbf{Ваш проект} & \textbf{Конкурент 1} & \textbf{Конкурент 2} \\ \hline
Цена & Средняя & Высокая & Низкая \\ \hline
Качество & Высокое & Среднее & Низкое \\ \hline
Сервис & 24/7 & 9-18 & Нет \\ \hline
\end{tabular}
\caption{Сравнительная таблица конкурентов}
\end{table}

% ==================== РАЗДЕЛ 4: ПРОДУКЦИЯ И УСЛУГИ ====================
\newpage
\section{Продукция и услуги}

\subsection{Описание продукта}
Подробно расскажите, что именно вы продаете. Какую проблему клиента это решает?

\subsection{Конкурентные преимущества}
Почему клиент выберет вас? (Уникальное торговое предложение - УТП).
\begin{enumerate}
    \item Уникальная технология.
    \item Высокий уровень сервиса.
    \item Экологичность материалов.
\end{enumerate}

% ==================== РАЗДЕЛ 5: МАРКЕТИНГОВАЯ СТРАТЕГИЯ ====================
\newpage
\section{Маркетинговая стратегия}

\subsection{Ценообразование}
Опишите стратегию цен: премиум, средний сегмент или эконом.

\subsection{Каналы продвижения}
Где и как вы будете рекламировать продукт?
\begin{itemize}
    \item Интернет-маркетинг (SEO, контекстная реклама).
    \item Социальные сети (SMM).
    \item Офлайн-мероприятия.
\end{itemize}

% ==================== РАЗДЕЛ 6: ФИНАНСОВЫЙ ПЛАН ====================
\newpage
\section{Финансовый план}

Это самая важная часть для инвестора. Все расчеты должны быть обоснованы.

\subsection{Стартовые затраты}
\begin{table}[h!]
\centering
\begin{tabular}{lc}
\toprule
\textbf{Статья расходов} & \textbf{Сумма (руб.)} \\
\midrule
Регистрация бизнеса & 50 000 \\
Закупка оборудования & 1 000 000 \\
Ремонт помещения & 500 000 \\
Маркетинг (старт) & 200 000 \\
\midrule
\textbf{Итого} & \textbf{1 750 000} \\
\bottomrule
\end{tabular}
\caption{Начальные инвестиции}
\end{table}

\subsection{Прогноз доходов}
\begin{table}[h!]
\centering
\begin{tabular}{lccc}
\toprule
\textbf{Показатель} & \textbf{1 год} & \textbf{2 год} & \textbf{3 год} \\
\midrule
Выручка & 5 000 000 & 8 000 000 & 12 000 000 \\
Расходы & 3 500 000 & 5 000 000 & 7 000 000 \\
\midrule
\textbf{Чистая прибыль} & \textbf{1 500 000} & \textbf{3 000 000} & \textbf{5 000 000} \\
\bottomrule
\end{tabular}
\caption{Прогноз прибылей и убытков}
\end{table}

\subsection{Окупаемость}
Рассчитайте срок окупаемости проекта (Payback Period). Приведите формулу расчета.

% ==================== ЗАКЛЮЧЕНИЕ ====================
\newpage
\section{Заключение}

Подведите итоги. Почему этот проект стоит инвестиций? Выразите уверенность в успехе и готовность предоставить дополнительные данные по запросу.

\end{document}
